% 第八章 催化剂与风险

\section{七大催化剂}

我们认为未来 12–18 个月内，功率半导体行业存在以下七大催化因素，可能成为板块行情的重要驱动：

\subsection{催化剂一：AI 资本开支超预期}

全球科技巨头 AI 资本开支持续超预期，英伟达、微软、Google、Meta 等公司的资本开支指引不断上调。AI 服务器需求旺盛直接拉动功率半导体需求，是行业最重要的边际催化。

\vspace{0.3em}

\textbf{关注指标}：北美五大科技公司资本开支指引、AI 服务器出货量、GPU 出货量。

\subsection{催化剂二：SiC 技术突破与成本快速下降}

SiC 衬底良率提升、8 英寸衬底量产、新生长工艺突破等技术进步有望加速 SiC 成本下降，推动渗透率快速提升。若 SiC 器件价格降至硅基 IGBT 的 1.5 倍以内，将迎来大规模替代窗口。

\vspace{0.3em}

\textbf{关注指标}：6 英寸衬底良率、8 英寸衬底量产进度、SiC 器件价格走势、主流车企 SiC 车型规划。

\subsection{催化剂三：车规级国产替代加速}

国内车规级 IGBT 和 SiC 器件正处于从验证到批量装车的关键期。若头部车企国产供应商份额快速提升，将显著提振市场对国产替代的信心。

\vspace{0.3em}

\textbf{关注指标}：国内车企 IGBT/SiC 国产化率、各公司车规业务营收占比、AEC-Q101 认证进展。

\subsection{催化剂四：行业价格企稳回升}

经过 2023–2024 年的价格调整，功率半导体产品价格已普遍回落 20–40\%，部分产品价格接近甚至低于现金成本。随着需求复苏和库存去化，价格有望企稳回升，带来业绩和估值的戴维斯双击。

\vspace{0.3em}

\textbf{关注指标}：渠道价格走势、晶圆代工产能利用率、厂商毛利率变化。

\subsection{催化剂五：光伏储能装机超预期}

全球光伏储能装机量持续超预期，拉动 IGBT 和 SiC 器件需求。尤其是高压大功率储能变流器对 SiC 的需求快速增长。

\vspace{0.3em}

\textbf{关注指标}：全球光伏新增装机、储能装机量、逆变器出货量。

\subsection{催化剂六：国产设备材料突破}

半导体设备和材料的国产化突破，将缓解国内功率半导体厂商的扩产瓶颈，降低对海外供应链的依赖，提升国产替代的长期确定性。

\vspace{0.3em}

\textbf{关注指标}：国产设备验证进度、关键材料国产化率、出口管制政策变化。

\subsection{催化剂七：政策支持力度加大}

半导体产业政策持续加码，大基金三期、税收优惠、补贴等政策支持有望超预期，为行业发展提供政策红利。

\vspace{0.3em}

\textbf{关注指标}：大基金投资进展、产业政策出台、税收优惠政策。

\section{七大风险因素}

在乐观的产业趋势背后，也需要警惕以下七大风险因素：

\begin{riskbox}[风险因素一览]
\small
\begin{tabularx}{\textwidth}{X L{3cm} c L{4.5cm}}
\toprule
\textbf{风险类别} & \textbf{风险事件} & \textbf{影响程度} & \textbf{关注信号} \\
\midrule
需求端风险 & AI 资本开支放缓 & \textcolor{riskred}{高} & 科技大厂CapEx指引下修、服务器出货不及预期 \\
需求端风险 & 新能源车销量不及预期 & \textcolor{riskred}{高} & 月度新能源车销量、车企排产计划 \\
供给端风险 & 产能过剩加剧 & \textcolor{riskamber}{中高} & 晶圆厂产能利用率、行业扩产计划 \\
技术迭代风险 & SiC/GaN 进展慢于预期 & \textcolor{riskamber}{中} & 良率提升进度、量产时间推迟 \\
\midrule
竞争格局风险 & 海外厂商降价竞争 & \textcolor{riskamber}{中} & 英飞凌等龙头价格策略变化 \\
政策风险 & 出口管制加码 & \textcolor{riskred}{高} & 美欧日半导体管制政策更新 \\
估值风险 & 板块估值回调 & \textcolor{riskamber}{中高} & 市场情绪变化、流动性收紧 \\
\bottomrule
\end{tabularx}
\end{riskbox}

以下对主要风险进行详细分析：

\subsection{风险一：AI 资本开支不及预期}

AI 是功率半导体最大的边际增量来源，如果 AI 投资周期见顶或增速放缓，将对行业成长性产生显著影响。需警惕的信号包括：科技大厂资本开支指引下修、GPU 需求放缓、AI 应用落地不及预期等。

\subsection{风险二：新能源车销量低于预期}

新能源汽车是功率半导体第一大下游，约占 35–40\% 的需求。若新能源车销量增速放缓或价格战加剧，将直接影响功率半导体需求。

\subsection{风险三：行业价格竞争加剧}

国内功率半导体厂商近年大规模扩产，如果需求增长不及供给扩张，可能引发价格战，压缩行业整体盈利能力。尤其需关注中低端 MOS、二极管等成熟产品的价格走势。

\subsection{风险四：SiC 技术迭代不及预期}

SiC 是行业重要的成长驱动，但如果 SiC 衬底良率提升慢于预期、成本下降不及预期，渗透率可能低于预测，影响相关公司的成长逻辑。

\subsection{风险五：国际贸易政策风险}

半导体是地缘政治博弈的核心领域，海外对中国半导体设备、材料、技术的出口管制可能进一步加码，影响国内厂商的扩产和技术升级。

\subsection{风险六：供给端扰动风险}

全球功率半导体供应链高度全球化，地缘冲突、自然灾害、疫情等不可预见因素可能对供应链造成扰动，带来短期供需失衡。

\subsection{风险七：估值波动风险}

功率半导体板块估值弹性大，市场情绪和流动性变化可能带来较大的估值波动。在高估值区间投资，需警惕估值回调风险。

\section{景气度跟踪指标}

为有效跟踪行业景气度变化，建议关注以下核心指标，按月度/季度频率进行跟踪验证：

\begin{exhibitbox}[表 10-1：功率半导体景气度跟踪指标体系]
\centering
\small
\begin{tabularx}{\textwidth}{X L{3cm} c L{4.5cm}}
\toprule
\textbf{指标类别} & \textbf{具体指标} & \textbf{更新频率} & \textbf{意义} \\
\midrule
\rowcolor{lightgrey}
\multicolumn{4}{l}{\textit{需求端指标}} \\
AI 服务器 & 服务器出货量 / 资本开支 & 季度 & 边际增量最大的下游 \\
新能源车 & 月度销量 / 排产计划 & 月度 & 第一大下游需求 \\
光伏储能 & 新增装机量 / 逆变器出货 & 月度 & 第三增长极 \\
工业 & PMI / 工控订单 & 月度 & 第二大下游 \\
\midrule
\rowcolor{lightgrey}
\multicolumn{4}{l}{\textit{供给端指标}} \\
晶圆代工 & 8英寸/12英寸产能利用率 & 季度 & 供给紧密度指标 \\
渠道库存 & 分销商库存周转天数 & 月度 & 渠道水位监测 \\
价格走势 & 主流型号价格指数 & 月度 & 景气度最直接反映 \\
\midrule
\rowcolor{lightgrey}
\multicolumn{4}{l}{\textit{上市公司验证指标}} \\
营收/毛利率 & 季度财报数据 & 季度 & 基本面验证 \\
存货变动 & 存货金额 / 周转天数 & 季度 & 库存周期验证 \\
资本开支 & 扩产计划 / 产能进度 & 半年度 & 未来供给预判 \\
\midrule
\rowcolor{lightgrey}
\multicolumn{4}{l}{\textit{技术进展指标}} \\
SiC 进度 & 衬底良率 / 客户认证 & 季度 / 事件驱动 & 技术迭代节奏 \\
GaN 进度 & 渗透率 / 新应用拓展 & 季度 / 事件驱动 & 新应用边界拓展 \\
车规认证 & AEC-Q101 认证进展 & 事件驱动 & 高端突破验证 \\
\bottomrule
\end{tabularx}
\sourcenote{本研究团队整理构建}
\end{exhibitbox}

\begin{keyinsight}[风险收益评估]
综合来看，我们认为功率半导体行业当前的风险收益比具有吸引力。上行空间来自 AI + 新能源双轮驱动的业绩增长和国产替代的份额提升，下行风险主要来自估值回调和行业周期波动。建议采取"核心 + 卫星"的配置策略，核心仓位配置低估值高确定性龙头，卫星仓位配置高成长高弹性标的，平衡收益与风险。
\end{keyinsight}
